\documentclass[12pt]{article}
\usepackage{amsmath, amssymb, amsthm}
\usepackage[margin=1in]{geometry}

\newtheorem{theorem}{Theorem}
\newtheorem{lemma}[theorem]{Lemma}

\title{Problem 234: Semidivisible Numbers}
\author{Project Euler Solution}
\date{}

\begin{document}
\maketitle

\section{Problem Statement}
For an integer $n$, let $\mathrm{lps}(n)$ be the largest prime $\leq \sqrt{n}$ and $\mathrm{ups}(n)$ the smallest prime $\geq \sqrt{n}$. A number is \emph{semidivisible} if exactly one of $\mathrm{lps}(n)$ and $\mathrm{ups}(n)$ divides it. Find the sum of all semidivisible numbers not exceeding $999\,966\,663\,333$.

\section{Mathematical Foundation}

\begin{theorem}[Interval Decomposition]
For consecutive primes $p < q$, every integer $n$ with $p^2 < n < q^2$ satisfies $\mathrm{lps}(n) = p$ and $\mathrm{ups}(n) = q$.
\end{theorem}

\begin{proof}
Since $p^2 < n < q^2$, we have $p < \sqrt{n} < q$. As $p$ and $q$ are consecutive primes, $p$ is the largest prime $\leq \sqrt{n}$ and $q$ is the smallest prime $\geq \sqrt{n}$.
\end{proof}

\begin{lemma}[Semidivisibility Criterion]
For $n \in (p^2, q^2)$ with consecutive primes $p < q$, $n$ is semidivisible iff exactly one of $p \mid n$, $q \mid n$ holds.
\end{lemma}

\begin{proof}
Immediate from the theorem and the definition of semidivisibility.
\end{proof}

\begin{lemma}[Arithmetic Sum Formula]
The sum of semidivisible numbers in the open interval $(p^2, \min(q^2, L))$ equals $\Sigma_p + \Sigma_q - 2\Sigma_{pq}$, where $\Sigma_d$ is the sum of multiples of $d$ in the interval.
\end{lemma}

\begin{proof}
Let $A$ (resp.\ $B$) be the set of multiples of $p$ (resp.\ $q$) in the interval. Semidivisible numbers form $A \triangle B$. Since $p,q$ are distinct primes, $A \cap B$ consists of multiples of $pq$. Then $\sum_{A \triangle B} = \sum_A + \sum_B - 2\sum_{A \cap B}$.
\end{proof}

\section{Algorithm}

\begin{verbatim}
function solve(L):
    primes = sieve(ceil(sqrt(L)) + margin)
    total = 0
    for each consecutive pair (p, q) in primes:
        lo = p*p, hi = min(q*q, L)
        if lo >= hi: continue
        total += sum_mult(p,lo,hi) + sum_mult(q,lo,hi) - 2*sum_mult(p*q,lo,hi)
    return total

function sum_mult(d, lo, hi):
    k1 = floor(lo/d) + 1, k2 = floor((hi-1)/d)
    if k1 > k2: return 0
    return d * (k1+k2) * (k2-k1+1) / 2
\end{verbatim}

\section{Complexity Analysis}
\begin{itemize}
    \item \textbf{Time:} $O(\sqrt{L}\,\log\log\sqrt{L})$ for the sieve; $O(\pi(\sqrt{L}))$ for the summation loop.
    \item \textbf{Space:} $O(\sqrt{L})$ for the sieve array.
\end{itemize}

\section{Answer}

\[
\boxed{1259187438574927161}
\]

\end{document}
